\chapter{La transition laminaire--turbulent}\label{ch:transition}
\iterbanner{V0}

La position de la transition est le paramètre qui pèse le plus sur la traînée d'un profil à $\Re\sim10^6$ : une couche limite turbulente frotte environ cinq fois plus qu'une couche laminaire. C'est aussi, on le verra, la principale source d'irrégularité de $C_D(m,p,t)$.

\section{Trois événements, le premier l'emporte}
Le code déclenche la transition à la première station où l'un des trois critères suivants est franchi.

\paragraph{1. Le critère de Michel (transition naturelle).} Corrélation empirique à un paramètre (Michel, ONERA, 1951) : la transition a lieu quand le Reynolds de quantité de mouvement dépasse une valeur qui dépend du Reynolds local,
\[
\Re_\theta>\Re_{\theta,tr}(\Re_x)=1{,}174\left(1+\frac{22\,400}{\Re_x}\right)\Re_x^{0{,}46},\qquad \Re_x=\frac{U_e\,s}{\nu}.
\]
Il représente la transition « naturelle » par amplification des ondes de Tollmien-Schlichting dans un environnement calme. Il ne contient ni le gradient de pression explicitement (seulement via l'historique dans $\theta$) ni le taux de turbulence extérieure.

\paragraph{2. Le décollement laminaire.} Si $\lambda<-0{,}09$ avant que Michel ne soit atteint, la couche laminaire décolle. Physiquement, elle forme le plus souvent un \textbf{bulbe de décollement} : la couche décollée transitionne très vite et recolle en turbulent. Le modèle suppose un bulbe court, réduit à une transition instantanée au point de décollement.

\paragraph{3. Le déclenchement forcé.} Une abscisse \code{x\_trip} imposée par l'utilisateur, qui représente une bande de rugosité. C'est la même sémantique que la commande \texttt{XTR} d'XFOIL : la transition a lieu au premier des deux, critère libre ou \code{x\_trip}. L'option \code{free\_transition=False} désactive Michel, pour pouvoir imposer une transition \emph{plus tardive} que celle de Michel, par exemple celle trouvée par XFOIL.

\begin{figure}[htbp]
\centering
% figure générée par tools/make_figures.py : transition_michel
\begin{tikzpicture}
\begin{groupplot}[group style={group size=2 by 1,horizontal sep=1.6cm},width=0.5\linewidth,height=5.4cm,
  xlabel={$x/c$},xmin=0,xmax=1,ymin=0,ymax=2000,title style={font=\small},tick label style={font=\scriptsize},
  legend style={draw=none,font=\scriptsize},legend pos=north west]
\nextgroupplot[title={extrados : transition en $x/c=0.279$ (michel)}]
\addplot[cblue,very thick] coordinates {(0.0038514,84.921) (0.0060621,102.74) (0.0087685,121.8) (0.011969,141.78) (0.01566,162.43) (0.01984,183.55) (0.024505,204.98) (0.029651,226.61) (0.035273,248.36) (0.041368,270.15) (0.047928,291.96) (0.054949,313.74) (0.062423,335.49) (0.070344,357.2) (0.078703,378.85) (0.087494,400.46) (0.096707,422.03) (0.10633,443.56) (0.11636,465.06) (0.12679,486.55) (0.13759,508.04) (0.14877,529.53) (0.16031,551.04) (0.17219,572.58) (0.18442,594.17) (0.19696,615.81) (0.20982,637.52) (0.22297,659.32) (0.2364,681.22) (0.25011,703.22) (0.26407,725.36) (0.27826,747.65) (0.29268,770.11) (0.30731,792.76) (0.32214,815.64) (0.33714,838.79) (0.3523,862.29) (0.36761,886.23) (0.38305,910.85) (0.39857,936.73) (0.41414,962.69) (0.42974,987.74) (0.44542,1012.1) (0.46115,1036.1) (0.47691,1059.9) (0.49269,1083.5) (0.50848,1106.9) (0.52425,1130.2) (0.54,1153.4) (0.5557,1176.6) (0.57135,1199.7) (0.58692,1222.7) (0.6024,1245.7) (0.61777,1268.7) (0.63302,1291.7) (0.64814,1314.8) (0.6631,1337.9) (0.6779,1361) (0.69252,1384.2) (0.70695,1407.6) (0.72116,1431.1) (0.73515,1454.8) (0.74891,1478.6) (0.76242,1502.8) (0.77566,1527.2) (0.78863,1551.9) (0.80131,1577) (0.81369,1602.5) (0.82576,1628.5) (0.8375,1655) (0.84891,1682.1) (0.85997,1709.9) (0.87068,1738.4) (0.88102,1767.8) (0.89098,1798.1) (0.90055,1829.5) (0.90973,1862) (0.9185,1895.9) (0.92686,1931.2) (0.9348,1968.2) (0.94232,2007.1) (0.9494,2048.2) (0.95603,2091.8) (0.96222,2138.4) (0.96795,2188.4) (0.97323,2242.6) (0.97804,2301.8) (0.98238,2367.2)};\addlegendentry{$Re_\theta$ (Thwaites)}
\addplot[corange,very thick] coordinates {(0.0038514,254.52) (0.0060621,268.16) (0.0087685,282.68) (0.011969,297.72) (0.01566,313.1) (0.01984,328.77) (0.024505,344.68) (0.029651,360.8) (0.035273,377.11) (0.041368,393.57) (0.047928,410.16) (0.054949,426.86) (0.062423,443.63) (0.070344,460.45) (0.078703,477.29) (0.087494,494.13) (0.096707,510.95) (0.10633,527.72) (0.11636,544.43) (0.12679,561.06) (0.13759,577.59) (0.14877,593.99) (0.16031,610.27) (0.17219,626.39) (0.18442,642.34) (0.19696,658.12) (0.20982,673.71) (0.22297,689.09) (0.2364,704.25) (0.25011,719.19) (0.26407,733.89) (0.27826,748.33) (0.29268,762.51) (0.30731,776.42) (0.32214,790.04) (0.33714,803.36) (0.3523,816.36) (0.36761,829.01) (0.38305,841.27) (0.39857,852.97) (0.41414,864.32) (0.42974,875.56) (0.44542,886.67) (0.46115,897.59) (0.47691,908.3) (0.49269,918.78) (0.50848,929.02) (0.52425,939.03) (0.54,948.78) (0.5557,958.29) (0.57135,967.53) (0.58692,976.51) (0.6024,985.23) (0.61777,993.68) (0.63302,1001.9) (0.64814,1009.8) (0.6631,1017.4) (0.6779,1024.7) (0.69252,1031.8) (0.70695,1038.5) (0.72116,1045) (0.73515,1051.2) (0.74891,1057.1) (0.76242,1062.6) (0.77566,1067.9) (0.78863,1072.8) (0.80131,1077.4) (0.81369,1081.7) (0.82576,1085.6) (0.8375,1089.2) (0.84891,1092.5) (0.85997,1095.4) (0.87068,1098) (0.88102,1100.2) (0.89098,1102) (0.90055,1103.4) (0.90973,1104.5) (0.9185,1105.1) (0.92686,1105.4) (0.9348,1105.2) (0.94232,1104.5) (0.9494,1103.5) (0.95603,1101.9) (0.96222,1099.8) (0.96795,1097.2) (0.97323,1094) (0.97804,1090.1) (0.98238,1085.5)};\addlegendentry{seuil de Michel}
\addplot[cmuted,dashed] coordinates {(0.2795,0) (0.2795,2000)};
\nextgroupplot[title={intrados : transition en $x/c=0.478$ (michel)}]
\addplot[cblue,very thick] coordinates {(0.0040239,28.372) (0.006247,39.504) (0.0089501,52.466) (0.01213,66.77) (0.015782,82.089) (0.019902,98.191) (0.024485,114.91) (0.029527,132.12) (0.035021,149.71) (0.040961,167.62) (0.047342,185.77) (0.054155,204.13) (0.061395,222.64) (0.069053,241.28) (0.077121,260.03) (0.085592,278.87) (0.094457,297.77) (0.10371,316.72) (0.11333,335.71) (0.12333,354.73) (0.13367,373.77) (0.14437,392.82) (0.1554,411.87) (0.16676,430.91) (0.17844,449.93) (0.19042,468.92) (0.20269,487.86) (0.21525,506.75) (0.22808,525.58) (0.24116,544.32) (0.25449,562.95) (0.26806,581.46) (0.28185,599.83) (0.29585,618.03) (0.31005,636.03) (0.32443,653.8) (0.33899,671.28) (0.3537,688.43) (0.36855,705.14) (0.38354,721.27) (0.39866,736.44) (0.41396,751.01) (0.42937,766.2) (0.44486,781.87) (0.4604,797.7) (0.47599,813.57) (0.4916,829.42) (0.50723,845.18) (0.52285,860.84) (0.53845,876.36) (0.55402,891.72) (0.56954,906.92) (0.58499,921.92) (0.60037,936.74) (0.61565,951.35) (0.63082,965.75) (0.64587,979.95) (0.66078,993.95) (0.67553,1007.7) (0.69012,1021.3) (0.70452,1034.7) (0.71872,1047.9) (0.73272,1060.9) (0.74649,1073.8) (0.76001,1086.5) (0.77329,1099.1) (0.7863,1111.6) (0.79903,1124) (0.81147,1136.3) (0.82361,1148.6) (0.83542,1160.9) (0.84691,1173.1) (0.85806,1185.5) (0.86886,1198) (0.8793,1210.6) (0.88936,1223.4) (0.89904,1236.4) (0.90832,1249.8) (0.9172,1263.6) (0.92567,1277.8) (0.93372,1292.7) (0.94134,1308.3) (0.94852,1324.7) (0.95526,1342.3) (0.96155,1361) (0.96737,1381.4) (0.97274,1403.7) (0.97763,1428.5) (0.98205,1456.4)};\addlegendentry{$Re_\theta$ (Thwaites)}
\addplot[corange,very thick] coordinates {(0.0040239,300.71) (0.006247,247.32) (0.0089501,235.02) (0.01213,237.11) (0.015782,245.52) (0.019902,257.09) (0.024485,270.39) (0.029527,284.69) (0.035021,299.6) (0.040961,314.88) (0.047342,330.38) (0.054155,346.02) (0.061395,361.72) (0.069053,377.43) (0.077121,393.12) (0.085592,408.77) (0.094457,424.33) (0.10371,439.81) (0.11333,455.18) (0.12333,470.44) (0.13367,485.56) (0.14437,500.55) (0.1554,515.39) (0.16676,530.08) (0.17844,544.61) (0.19042,558.98) (0.20269,573.19) (0.21525,587.23) (0.22808,601.11) (0.24116,614.82) (0.25449,628.36) (0.26806,641.73) (0.28185,654.94) (0.29585,667.99) (0.31005,680.89) (0.32443,693.63) (0.33899,706.23) (0.3537,718.71) (0.36855,731.07) (0.38354,743.36) (0.39866,755.68) (0.41396,768.01) (0.42937,780.01) (0.44486,791.69) (0.4604,803.1) (0.47599,814.26) (0.4916,825.19) (0.50723,835.9) (0.52285,846.39) (0.53845,856.66) (0.55402,866.71) (0.56954,876.55) (0.58499,886.18) (0.60037,895.59) (0.61565,904.79) (0.63082,913.77) (0.64587,922.54) (0.66078,931.08) (0.67553,939.41) (0.69012,947.51) (0.70452,955.38) (0.71872,963.02) (0.73272,970.43) (0.74649,977.61) (0.76001,984.53) (0.77329,991.21) (0.7863,997.64) (0.79903,1003.8) (0.81147,1009.7) (0.82361,1015.4) (0.83542,1020.7) (0.84691,1025.8) (0.85806,1030.6) (0.86886,1035) (0.8793,1039.2) (0.88936,1043) (0.89904,1046.6) (0.90832,1049.7) (0.9172,1052.6) (0.92567,1055) (0.93372,1057.1) (0.94134,1058.7) (0.94852,1060) (0.95526,1060.8) (0.96155,1061.1) (0.96737,1060.8) (0.97274,1060) (0.97763,1058.6) (0.98205,1056.4)};\addlegendentry{seuil de Michel}
\addplot[cmuted,dashed] coordinates {(0.4782,0) (0.4782,2000)};
\end{groupplot}
\end{tikzpicture}

\caption{NACA 2412, $\alpha=2^\circ$, $\Re=3\cdot10^6$ : $\Re_\theta$ laminaire calculé par Thwaites (bleu) et seuil de Michel (orange) le long de chaque face. La transition a lieu au croisement, interpolé entre les deux stations qui l'encadrent.}\label{fig:michel}
\end{figure}

\section{Interpoler la position de transition}
Chaque critère s'écrit comme une fonction $g_i$ aux stations, avec la convention « franchi si $g\ge0$ » :
\[
g^{\text{Michel}}=\Re_\theta-\Re_{\theta,tr}(\Re_x),\qquad g^{\text{décollement}}=-0{,}09-\lambda,\qquad g^{\text{forcé}}=x-x_{trip}.
\]
À la première station $i$ où l'un d'eux devient positif, on interpole linéairement le zéro entre $i-1$ et $i$ :
\[
w=\frac{-g_{i-1}}{g_i-g_{i-1}}\in[0,1],\qquad s_{tr}=s_{i-1}+w\,(s_i-s_{i-1}).
\]
La condition initiale de Head (ch.~\ref{ch:head}) est prise en $s_{tr}$, avec $\theta$ interpolé au même endroit.

\paragraph{Pourquoi c'est important.} Sans interpolation, $x_{tr}$ ne peut prendre que les abscisses des stations. Quand $(m,p,t)$ varie continûment, $x_{tr}$ saute d'une station à la suivante et $C_D$ devient une fonction en escalier. Mesuré sur le NACA 2412, $C_D$ variait de $\pm1{,}5$~\% quand on changeait le nombre de panneaux. Avec l'interpolation, il varie de $\pm0{,}1$~\%. Pour un ajustement quadratique, des marches sont du bruit pur.

\begin{codebox}[ : recherche de la transition dans \code{march\_surface}]
\begin{lstlisting}
crit = {}
if x_trip is not None:
    crit["forced"] = x - x_trip
crit["laminar_sep"] = np.where(s > 0.002, LAMBDA_SEP - lam, -1.0)
if free_transition:
    crit["michel"] = Re_th - michel_transition(Re_x)
i_tr, mode, w = n, "none", 0.0
for i in range(n):
    hit = [k for k, g in crit.items() if g[i] >= 0]
    if hit:
        i_tr, mode = i, hit[0]
        if i > 0:
            g0, g1 = crit[mode][i - 1], crit[mode][i]
            w = float(np.clip(-g0 / (g1 - g0), 0.0, 1.0)) if g1 != g0 else 1.0
        break
\end{lstlisting}
Le décollement laminaire est ignoré sur les deux premiers millièmes de corde ($s\le0{,}002$) : juste après le point d'arrêt, $U_e'$ est mal estimé par différences finies et $\lambda$ peut être faussement très négatif. L'ordre des clés du dictionnaire fixe la priorité en cas d'égalité sur une même station : forcé, puis décollement laminaire, puis Michel.
\end{codebox}

\section{Comportement sur un profil}
\begin{figure}[htbp]
\centering
% figure générée par tools/make_figures.py : xtr_alpha
\begin{tikzpicture}
\begin{axis}[width=0.7\linewidth,height=5.6cm,xlabel={$\alpha$ (°)},ylabel={$x_{tr}/c$},xmin=-4,xmax=10,ymin=0,ymax=1,
  legend style={draw=none,font=\scriptsize,at={(0.02,0.02)},anchor=south west}]
\addplot[cblue,very thick,mark=*,mark size=1.4pt] coordinates {(-4,0.66149) (-3,0.59365) (-2,0.52539) (-1,0.46) (0,0.40359) (1,0.34723) (2,0.27945) (3,0.20801) (4,0.14399) (5,0.098468) (6,0.070947) (7,0.054723) (8,0.022071) (9,0.013933) (10,0.010145)};\addlegendentry{extrados}
\addplot[corange,very thick,mark=square*,mark size=1.4pt] coordinates {(-4,0.072542) (-3,0.096003) (-2,0.13131) (-1,0.18121) (0,0.24686) (1,0.33239) (2,0.47818) (3,0.61065) (4,0.7494) (5,0.8636) (6,0.9027) (7,0.91429) (8,0.92399) (9,0.93213) (10,0.93901)};\addlegendentry{intrados}
\addplot[only marks,mark=o,mark size=3.5pt,cink,thick] coordinates {(6,0.9027) (7,0.91429) (8,0.92399) (9,0.93213) (10,0.93901) (8,0.022071) (9,0.013933) (10,0.010145)};\addlegendentry{décollement laminaire}
\end{axis}
\end{tikzpicture}

\caption{Position de transition du NACA 2412 en fonction de l'incidence ($\Re=3\cdot10^6$). Quand $\alpha$ augmente, la transition remonte vers le bord d'attaque sur l'extrados (gradient adverse plus fort après le pic d'aspiration) et recule sur l'intrados. Cercles : transition par décollement laminaire plutôt que par Michel.}\label{fig:xtr-alpha}
\end{figure}

La figure~\ref{fig:xtr-alpha} montre le comportement attendu, et une bascule de mode. Sur l'intrados, à partir de $\alpha=6^\circ$, la couche laminaire décolle avant d'atteindre le seuil de Michel. Sur l'extrados, la même bascule se produit vers $\alpha=8^\circ$, sous le pic d'aspiration. Chaque bascule crée une rupture de pente dans $x_{tr}(\alpha)$, donc dans $C_D(\alpha)$.

\begin{figure}[htbp]
\centering
% figure générée par tools/make_figures.py : cd_xtrip
\begin{tikzpicture}
\begin{axis}[width=0.7\linewidth,height=5.4cm,xlabel={transition imposée $x_{tr}/c$ (deux faces)},ylabel={$C_D$ (counts)},
  xmin=0,xmax=0.8,legend style={draw=none,font=\scriptsize},legend pos=north east]
\addplot[cblue,very thick,mark=*,mark size=1.2pt] coordinates {(0.02,97.593) (0.05,95.511) (0.08,93.023) (0.11,90.368) (0.14,87.62) (0.17,84.83) (0.2,82.031) (0.23,79.231) (0.26,76.435) (0.29,73.664) (0.32,70.911) (0.35,68.203) (0.38,65.532) (0.41,63.877) (0.44,63.043) (0.47,62.211) (0.5,61.374) (0.53,60.543) (0.56,59.706) (0.59,58.873) (0.62,58.04) (0.65,57.206) (0.68,56.373) (0.71,55.538) (0.74,54.704) (0.77,53.87) (0.8,53.035)};\addlegendentry{transition imposée}
\addplot[corange,dashed,very thick] coordinates {(0,69.17) (0.8,69.17)};\addlegendentry{transition libre (Michel)}
\end{axis}
\end{tikzpicture}

\caption{$C_D$ du NACA 2412 ($\alpha=2^\circ$) quand on impose la même position de transition sur les deux faces. De $x_{tr}=0{,}02$ à $0{,}8$, la traînée passe de 98 à 53~counts. En transition libre, Michel donne 69~counts.}\label{fig:cd-xtrip}
\end{figure}

La figure~\ref{fig:cd-xtrip} chiffre l'enjeu : selon l'endroit où la transition a lieu, la traînée du même profil varie presque du simple au double. Toute erreur sur le critère de transition se retrouve directement dans $C_D$.

\begin{piege}
Michel est connu pour être \textbf{trop précoce à bas Reynolds}. Sur le NACA 0012 à $\alpha=0$ et $\Re=10^6$, il fait transitionner en $x/c=0{,}46$, et donne $C_D=75$~counts contre environ 54 couramment cités pour XFOIL, à reconfirmer au test V5. Si l'on impose la transition au décollement laminaire ($x\approx0{,}62$), on obtient 62~counts. L'erreur de transition domine donc l'erreur de couche limite. La méthode $e^N$ d'XFOIL suit l'amplification des ondes instables et voit l'effet stabilisant d'un gradient favorable ; Michel ne le voit pas.
\end{piege}

\begin{projet}
La bascule entre transition de Michel et transition par décollement laminaire est la principale non-régularité de $C_D(m,p,t)$ sur la grille. C'est elle qui limite le $R^2$ du modèle quadratique à environ 0,91 (ch.~\ref{ch:espace}). Deux leviers existent si la phase~4 en a besoin. On peut forcer la transition près du bord d'attaque (cas « tout turbulent », plus régulier et défendable pour un drone de série aux surfaces imparfaites). On peut aussi caler le critère sur XFOIL.
\end{projet}

\exercices
\begin{exercice}
Sur plaque plane, Thwaites donne $\Re_\theta=0{,}671\sqrt{\Re_x}$. Écrire l'équation qui donne le $\Re_x$ de transition selon Michel, et la résoudre numériquement (quelques itérations suffisent). En déduire $x_{tr}$ pour $\Re=3\cdot10^6$.
\end{exercice}
\begin{exercice}
On fait varier continûment l'épaisseur $t$ d'un profil. Expliquer pourquoi, sans interpolation de $x_{tr}$, la courbe $C_D(t)$ présenterait de petites marches. Quel serait l'effet de ces marches sur un modèle quadratique ajusté par moindres carrés, et sur un optimiseur par gradient ?
\end{exercice}
